\section{Related Work}

\aawarning{TODO---use~\citep{santoroSimpleNeuralNetwork2017,santoroRelationalRecurrentNeural2018,webbEmergentSymbolsBinding2021} for relevant OG references.}

\begin{itemize}
    \item The ability for neural networks to learn to perform relational reasoning in a data-efficient manner while achieving generalization through the right types of abstractions.
    \item This is an ability humans possess, but which neural models struggle with (in terms of efficiency).
    \item One theory, which can be described in the language of ``inductive biases'', that has emerged for this pertains to the distinction between so-called ``sensory'' and ``relational'' information. Modern neural architectures generally lack explicit computational mechanisms for computing feature representations that explicitly encode relational information (that is, comparisons). Although neural models can learn to fit their data on relational tasks given enough data, this lack of relational inductive bias results in limited generalization and requires large amounts of data.
    \item More recently, there has been a significant amount of work evaluating the relational reasoning abilities of large language models based on the Transformer architecture. This work finds that some level of relational reasoning ability is attained in large models after injesting large amounts of data. Moreover, there is evidence that the internal representations come to encode interpretable relational representations.
    \item In response, a recent line of work studied the problem of designing neural architectures for enhanced relational reasoning, seeking to embed relational inductive biases which drive the model to discover solutions which generalize better with less data.
    \item An early example is~\citep{santoroSimpleNeuralNetwork2017} which ... .
    \item More examples...
    \item The most related work to ours is ~\citep{altabaa2024abstractors} which proposes an attention mechanism where the values are replaced by ..., 
    \item Previous work has been more restricted in domain, and mainly applied to synthetic relational benchmarks. The goal of the present work is ...
    \item Relationship to symbolic approaches to AI, which are inherently relational.
\end{itemize}